% !TEX TS-program = XeLaTeX
% use the following command:
% all document files must be coded in UTF-8
\documentclass[english]{textolivre}
% build HTML with: make4ht -e build.lua -c textolivre.cfg -x -u article "fn-in,svg,pic-align"

\journalname{Texto Livre}
\thevolume{19}
%\thenumber{1} % old template
\theyear{2026}
\receiveddate{\DTMdisplaydate{2026}{3}{30}{-1}} % YYYY MM DD
\accepteddate{\DTMdisplaydate{2026}{4}{30}{-1}}
\publisheddate{\DTMdisplaydate{2026}{6}{18}{-1}}
\corrauthor{Francis Arthuso Paiva}
\articledoi{10.1590/1983-3652.2026.65345}
%\articleid{NNNN} % if the article ID is not the last 5 numbers of its DOI, provide it using \articleid{} commmand 
% list of available sesscions in the journal: articles, dossier, reports, essays, reviews, interviews, editorial
\articlesessionname{dossier}
\runningauthor{Kalantzis et al.} 
%\editorname{Leonardo Araújo} % old template
\sectioneditorname{Daniervelin Pereira~\orcid{0000-0003-1861-3609}}
\layouteditorname{Leonardo Araujo~\orcid{0000-0003-3884-2177}}

\title{Interviews and talks with Bill Cope, Mary Kalantzis, and Vânia Castro}
\othertitle{Entrevistas e conversas com Bill Cope, Mary Kalantzis e Vânia Castro}
% if there is a third language title, add here:
%\othertitle{Artikelvorlage zur Einreichung beim Texto Livre Journal}


\author[1]{Mary Kalantzis~\orcid{0000-0002-6552-3780}\thanks{Email: \href{mailto:kalantzi@illinois.edu}{kalantzi@illinois.edu}}}
\author[1]{Bill Cope~\orcid{0000-0003-2129-395X}\thanks{Email: \href{mailto:billcope@illinois.edu}{billcope@illinois.edu}}}
\author[1]{Vania Carvalho de Castro~\orcid{0000-0002-4930-5345}\thanks{Email: \href{mailto:vaniac@illinois.edu}{vaniac@illinois.edu}}}
\author[2]{Záira Bomfante dos Santos~\orcid{0000-0002-6162-8489}\thanks{Email: \href{mailto:zbomfante@gmail.com}{zbomfante@gmail.com}}}
\author[3]{Francis Arthuso Paiva~\orcid{0000-0002-9083-3342}\thanks{Email: \href{mailto:paivafrancis@yahoo.com.br}{paivafrancis@yahoo.com.br}}}
\author[4]{Fei Victor Lim~\orcid{0000-0003-3046-1011}\thanks{Email: \href{mailto:victor.lim@nie.edu.sg}{victor.lim@nie.edu.sg}}}
\author[5]{Camila Cárdenas Neira~\orcid{0000-0002-1842-7200}\thanks{Email: \href{mailto:camila.cardenas.neira@gmail.com}{camila.cardenas.neira@gmail.com}}}
\affil[1]{University of Illinois, Urbana-Champaign, Illinois, USA.}
\affil[2]{Universidade Federal do Espírito Santo, Vitória, ES, Brasil.}
\affil[3]{Universidade Federal de Minas Gerais, Colégio Técnico, Belo Horizonte, MG, Brasil.}
\affil[4]{Nanyang Technological University, Nanyang Ave, Singapore.}
\affil[5]{Universidad Austral de Chile, Valdivia, Los Ríos, Chile.}


\addbibresource{article.bib}
% use biber instead of bibtex
% $ biber article

% used to create dummy text for the template file
\definecolor{dark-gray}{gray}{0.35} % color used to display dummy texts
\usepackage{lipsum}
\SetLipsumParListSurrounders{\colorlet{oldcolor}{.}\color{dark-gray}}{\color{oldcolor}}

% used here only to provide the XeLaTeX and BibTeX logos
\usepackage{hologo}

% if you use multirows in a table, include the multirow package
\usepackage{multirow}

% provides sidewaysfigure environment
\usepackage{rotating}

% CUSTOM EPIGRAPH - BEGIN 
%%% https://tex.stackexchange.com/questions/193178/specific-epigraph-style
\usepackage{epigraph}
\renewcommand\textflush{flushright}
\makeatletter
\newlength\epitextskip
\pretocmd{\@epitext}{\em}{}{}
\apptocmd{\@epitext}{\em}{}{}
\patchcmd{\epigraph}{\@epitext{#1}\\}{\@epitext{#1}\\[\epitextskip]}{}{}
\makeatother
\setlength\epigraphrule{0pt}
\setlength\epitextskip{0.5ex}
\setlength\epigraphwidth{.7\textwidth}
% CUSTOM EPIGRAPH - END

% to use IPA symbols in unicode add
%\usepackage{fontspec}
%\newfontfamily\ipafont{CMU Serif}
%\newcommand{\ipa}[1]{{\ipafont #1}}
% and in the text you may use the \ipa{...} command passing the symbols in unicode

% LANGUAGE - BEGIN
% ARABIC
% for languages that use special fonts, you must provide the typeface that will be used
% \setotherlanguage{arabic}
% \newfontfamily\arabicfont[Script=Arabic]{Amiri}
% \newfontfamily\arabicfontsf[Script=Arabic]{Amiri}
% \newfontfamily\arabicfonttt[Script=Arabic]{Amiri}
%
% in the article, to add arabic text use: \textlang{arabic}{ ... }
%
% RUSSIAN
% for russian text we also need to define fonts with support for Cyrillic script
% \usepackage{fontspec}
% \setotherlanguage{russian}
% \newfontfamily\cyrillicfont{Times New Roman}
% \newfontfamily\cyrillicfontsf{Times New Roman}[Script=Cyrillic]
% \newfontfamily\cyrillicfonttt{Times New Roman}[Script=Cyrillic]
%
% in the text use \begin{russian} ... \end{russian}
% LANGUAGE - END

% EMOJIS - BEGIN
% to use emoticons in your manuscript
% https://stackoverflow.com/questions/190145/how-to-insert-emoticons-in-latex/57076064
% using font Symbola, which has full support
% the font may be downloaded at:
% https://dn-works.com/ufas/
% add to preamble:
% \newfontfamily\Symbola{Symbola}
% in the text use:
% {\Symbola }
% EMOJIS - END

% LABEL REFERENCE TO DESCRIPTIVE LIST - BEGIN
% reference itens in a descriptive list using their labels instead of numbers
% insert the code below in the preambule:
%\makeatletter
%\let\orgdescriptionlabel\descriptionlabel
%\renewcommand*{\descriptionlabel}[1]{%
%  \let\orglabel\label
%  \let\label\@gobble
%  \phantomsection
%  \edef\@currentlabel{#1\unskip}%
%  \let\label\orglabel
%  \orgdescriptionlabel{#1}%
%}
%\makeatother
%
% in your document, use as illustraded here:
%\begin{description}
%  \item[first\label{itm1}] this is only an example;
%  % ...  add more items
%\end{description}
% LABEL REFERENCE TO DESCRIPTIVE LIST - END


% add line numbers for submission
%\usepackage{lineno}
%\linenumbers


\begin{document}
\maketitle

\begin{polyabstract}
\begin{abstract}
Interview with Bill Cope, Mary Kalantzis, and Vânia Castro on Artificial Intelligence (AI), teaching and learning from the perspective of design-based learning, multiliteracy pedagogy and multimodality, conducted online in December 2025. The questions were based on the aforementioned professors’ articles about their proposal for transpositional grammar and their AI environment called Cyber Scholar. The interview consists of five thematic questions with their respective answers and references. Its content is consistent with the special issue proposal.

\keywords{Artificial Intelligence \sep Multiliteracy pedagogy \sep Multimodality}
\end{abstract}

\begin{english}
\begin{abstract}
Entrevista com Bill Cope, Mary Kalantzis e Vânia Castro sobre Inteligência Artificial (IA), ensino e aprendizagem na perspectiva da aprendizagem por design, da pedagogia dos multiletramentos e da multimodalidade, realizada virtualmente em dezembro de 2025. As questões foram produzidas com base nos artigos dos referidos professores acerca da sua proposta de transpositional grammar e do seu ambiente de IA chamado Cyber Scholar. A entrevista consiste em cinco perguntas temáticas com suas respectivas respostas e referências. Seu conteúdo está em consonância com a proposta do dossiê.

\keywords{Inteligência Artificial \sep Pedagogia dos multiletramentos \sep Multimodalidade}
\end{abstract}
\end{english}
% if there is another abstract, insert it here using the same scheme
\end{polyabstract}

\section{Question 1}\label{sec-q1}

% Interview Begin
\setinterviewee{Mary Kalantzis and Bill Cope}{}{University of Illinois, Urbana- Champaign, Illinois, USA}{billcope@illinois.edu}
\setinterviewer{Záira Bomfante dos Santos}{}{Universidade Federal do Espírito Santo, Vitória, ES, Brasil}{zbomfante@gmail.com}
\renewcommand{\interviewdate}{\today}
\renewcommand{\interviewcontext}{
    \textbf{Context:} This is a template interview.
}
\setspeaker{Qu}[\interviewerfname]
\setspeaker{An}[\intervieweefname]
\renewcommand{\interviewcontent}{
    \Qu The first question is related to the notion of grammar. In your recent published works, \textcite{cope2024,kalantzis2025}, two terms grabbed my attention: multimodal grammar of artificial intelligence and cyber-social literacy learning. Considering the notion of grammar involved in these works, it is anchored in Halliday’s view that a grammar is a resource for meaning, and that literacy from this perspective no longer affords to be instrumental and functional. In the age of AI, how can we develop literacy that leads students to a critical awareness of the patterns of meaning that are produced? How can we establish an active relationship with the machine as a result of interested human agency?
    \An The first question is related to the notion of grammar. I want to make a statement before we start, because your questions are very profound, and I just want to put it up in your context. I just want to say to you that Bill and I have been theorists and scholars. But that’s not all we’ve been. \\
    From the beginning of our careers, and this is important to us, we have also been practitioners. So we’ve always taken theory and tried to put it into practice. We started very early with genre theory and what we call social literacy, where we wanted people to understand that meaning-making had a purpose. It was important for everybody to know how they became who they were and how their choices influenced others. So we created teaching materials for primary school and high school. In Australia, they were used everywhere. Then, after that, we did learning by design, which again, theoretically, was about pedagogy and how educators had to be purposeful and understand every student, every class, every content goal. We produced a material and guidelines about how to do learning by design, which worked across Australia and outside Australia. Since then, we moved to multiliteracies because we thought two things mattered: multimodality in the digital age and multilingualism in all its forms. Thus, with those two things, diversity in meaning-making in languages and the affordances of the new technology, we put together models about how to interpret that in practice. And that was followed by new learning and other work that we put together to say, we are in an epochal shift, and we need to carry on with shifts.\\
    The new learning book was translated into Portuguese, Greek, and many other languages, which was a call to say we can’t go on doing what we’ve done. We have to rethink it. Of course, we put out the transpositional grammar because we thought that traditional alphabetical literacy, which is labels, right, cannot understand this epoch, which just shifts. We needed a grammar that saw how shifts were made between the different modes and the different functions. Now we’ve arrived at another moment, AI and the digital taking us even further, and again, what we’re trying to do is to say, well, what does it look like? So, we as educators can’t just write textbooks, can’t just produce instructional design. We have to create new tools in collaboration with others. So at every point, we move backwards and forwards from theory to practice to testing. Thereby, we just wanted to say that at the beginning because this moment is another epochal shift, which has tremendous implications for all teaching and all learning. All of us are not quite sure how to deal with it at the moment, but we wanted you to see continuously the kind of scholars we’re being, not just theoreticians, you know, but also testing how that theory works in practice. So, we just wanted to share that with you at the beginning.
    \Qu We greatly appreciate your statement. Please, tell us more about your notion of transpositional grammar.
    \An Grammar represents patterns of meaning. How is meaning patterned? And traditionally, you know, we’ve had theories that come from Chomsky, for example, that are about syntax, mostly, and the structure of language. And then, we have Halliday, who moves towards social functions, right, and is not so abstract as Chomsky. But we push it further, right? One, we push the idea of grammar further into multimodality, not just alphabetical or symbolic literacy, and two, we move beyond the sentence and the lexicon in order to go to meaning as context, meaning as interest, meaning influenced by the social context in which they’re made, because all meaning starts with a sense of purpose. And it’s not just narrow linguistics, not just, you know, traditional grammar. You’ll have to understand how the purpose shapes the  grammar. \\
    If you’ve got somebody there who can actually combine these two things, one of the things is what Chomsky did, so this is kind of understanding the whole of linguistics in a way, and I was saying, where are we sitting in relation to that linguistic division? So what Chomsky did is: he had a theory of syntax, and what Chomsky said is we’re not dealing with meaning. We’re only dealing with language as the form, the way in which words are assembled, not like society. We’re just dealing with the fact that there are these abstract things, nouns and verbs, and we put them together, and that’s grammar. So in a sense, what we’re talking about here is the definition of grammar, what we mean by grammar; that’s what he meant by grammar. One of the things Chomsky says quite explicitly is he never meant to be talking about meaning. On the other hand, he also never meant any of this theory to be of any use in education, and it hasn’t been.\\
    Thus, we move on now to Halliday. What Halliday does is talk about meaning, because he says that the way to understand grammar is by its functions. So again, if we go back to that sentence with nouns and verbs, you know, nouns have a certain function, which in our new, our more recent grammar, is a referencing of the world and native things in the world, and verbs have a certain function, which is that they describe certain kinds of action. In other words, what we’ve got there is we’ve got kind of a little theory of the world, actually, which consists of things and actions. So, what Halliday does is call this a functional grammar, because it’s about things that you do, but what he’s actually doing is trending towards semantics. He’s actually talking about kind of structures of meaning. Now, one of the things we do with our grammar is what we call transpositional grammar, which is our most recent iteration of this.\\
    We put it beyond Halliday in two ways. One way is we deal with context as well, so we actually have that as part of our grammar. In fact, Michael Halliday’s wife retired her son, worked on that, and it seems that Mary was taught by her. One of the things was, and this is actually the tradition of the domain of pragmatics, that you can’t understand meaning by just watching a sentence. You have to understand the context in which it’s uttered and the things to which it’s referring. So, we’ve kind of built a theory of context into what we’re doing as well, but one of the other key concepts behind the idea of a function is purpose. It’s very hard to understand why a meaning takes its particular form and how it’s organized, except by its differential purposes. An advertisement might talk about a car, which is different from the way you and we talk about a car. So it’s all cars, but their purposes explain discourse differences between an advertisement and a conversation. That’s a simple example. So one of the things is the function then becomes a very important idea.\\
    In our transpositional grammar, we have a whole area called interest, which tries to theorize that. But in a way, both context and interest are actually outside the Hallidayan system. Our argument is you need to bring it into the system. In our system, we’ve got the three Hallidayan metafunctions, and we’ve renamed them as reference, agency, and structure. We mean something slightly different, but it’s more or less the same and we’ve added these two other components, context and interest, which expand the Hallidayan scheme. But the other thing is the Hallidayan scheme was really only about language. One of the things is that language barely, rarely exists by itself anymore. This is the whole realm of multimodality. So we’ve got a kind of theory about these multiple forms of meaning, and we wanted to build a grammar that described all those forms of meaning, which are not just text, but image and space and object and body and sound and speech, all of those things. We wanted to have kind of a mental level grammar that described all those things. For example, when you come to reference, you might use a noun, or when you come to an image, it’s a picture, it’s the same thing. Here’s the word car, here’s a picture of a car; it’s the same thing, but obviously each of those forms of meaning has different importance.\\
    So, what we want to do is, rather than use the word noun, we want to talk about reference. The difference between reference, let’s say, a common noun or a proper noun is that a common noun is cars, and a proper noun is this particular car with this registration number. The registration number is actually a proper noun. We wanted to actually have this idea of these high-level functions. Therefore, what we then do is, instead of within reference, instead of common nouns and proper nouns, we have instances and concepts, and then we can use that word to apply across all the forms of meaning. There are basically two dimensions where we’ve extended Halliday. One is the functions, adding two more functions, and then building a grammar that goes across all these forms.\\
    Now, coming back to your question about grammar, this means that we’re not dealing with grammar in either the same sense as Chomsky, who said it was not semantics, or Halliday, who used grammar in a linguistic kind of way. We reach the point where we’re using grammar as a kind of metaphor, and it’s a metaphor for parsing, to use another language metaphor, parsing the patterns of meaning there are in the world. Moreover, this is because multimodality really matters as the multiple functions. You needed a grammar that went across the modes, not a grammar tied just to alphabetical literacy, and that valued each of those modes equally and could parse them. You could parse them even when they were mixed. So, that took us to the transpositional grammar. It’s important for us, but it’s a very hard, very difficult shift to make because traditional alphabetical grammar is what everyone learns. Yet, meaning-making to us is always moving between writing, speaking, visual, and audio, and we just wanted to value all the modes equally and have a means of parsing them with equal values. For example, the words on a TikTok video. The interesting TikTok video is going to have words streaming over it all the time. The words on a TikTok video make no sense or make little sense except in relation to the image and the embodied gestures of people. So, there are a lot of languages in this era, this kind of inequity. But let me just mention one other key concept in this transfer of grammar, that’s the word transposition. In addition, what we are doing is we are shifting a meaning from one space to another. We’ve got a picture of the car and we’ve got the word car; what we can do is we can move backwards and forwards between these. That’s transposition, which is moving across on that dimension. But on the function dimension, we’re moving all the time between one thing and the other. So, for example, let’s take Mary’s name. Mary is a person, a common noun, and Mary is a proper noun. Well, what we do is recognize that those meanings are simultaneous. We move between them interchangeably. Mary is both a person, both a common noun and a proper noun.\\
    Thus, what we do is we talk about these kinds of movements backwards and forwards as transposition. So functional movement on one dimension and then this multimodal with different forms of meaning on the other dimension. That’s what we mean by grammar essentially. And the multimodal. But even in the textual world, you’ve got these other movements along the other dimension.
    \Qu I think it’s crucial for us in Brazil to break this idea of grammar as something rigid, something related only to verbal mode. It’s so crucial for us and the educational system in Brazil.
    \An The grammar that was always there is still there, but it’s just there in text, and what we want to do is we want to expand that. So we are not saying that the grammar that Chomsky and Halliday found in text is not there. Of course it’s there, but we just want to kind of expand the grammar to the realities of what’s happened post the digital age.
}
    \begin{description}
        \interviewcontent
    \end{description}
% Interview End




\section{Question 2}\label{sec-q2}
% Interview Begin
\setinterviewee{Mary Kalantzis and Bill Cope}{}{University of Illinois, Urbana- Champaign, Illinois, USA}{john@example.com}
\setinterviewer{Fei Victor Lim}{}{Nanyang Technological University, Nanyang Ave, Singapore}{victor.lim@nie.edu.sg}
\renewcommand{\interviewdate}{\today}
\renewcommand{\interviewcontext}{
    \textbf{Context:} This is a template interview.
}
\setspeaker{Qu}[\interviewerfname]
\setspeaker{An}[\intervieweefname]
\renewcommand{\interviewcontent}{
    \Qu You have discussed both the promises and perils of using Generative AI (GenAI) in teaching and learning in your papers. What do you see as the key design principles that should underpin the development of GenAI tools for educational contexts?
    \An The first thing we have to say is that AI did not come out of the academy. It did not come from us. It came from six men in America, four in China, one in France, who scraped everything that’s available to them on the internet, everything, and now give it back to us for a fee. They’ve commercialized it; you pay \$20 or \$200. First of all, we have to recognize that, at the moment, it’s benign. We don’t know where it’s going to go, but they control the way in which we engage with it. So, people are correct to be anxious. They control it at every level; besides, they fine-tune the way the chat box relates to you. They have to be polite, nice, and kind of do what you want. That’s not in the corpus, that’s in the fine-tuning. The corpus is there with everything. Also, there’s fine-tuning how you engage, how we engage with the corpus that is created by these owners. At the moment, in America, there’s no regulation of them. Europe is trying to regulate them, but they are controlled by those owners who have commercial interests. They want us engaged all the time. All the time engaged. They find means of doing that for us so that they can then monetize that engagement in some way. Thus, first of all, we have to understand that. \\
    I (Bill) want to say that this is very different from the foundation of the Internet. The Internet is pretty open and it’s pretty free, but the great irony is that, in the beginning, the Internet was funded by the U.S. military. So, U.S. universities built the Internet and for this it became open and it became free. We could be charging now for every email you sent, but you sent emails for nothing. You sent emails for free. We could be charging them every time you use the Internet to access the website. We don’t. The irony is it was public funding in the U.S., the U.S. military and universities, which created a relatively open Internet. The foundation models are not open. They’ve been created by private companies on massive capital investments. In fact, now we’re talking about an AI bubble where the investments might have been vastly, vastly overplayed. They’re talking about hundreds of billions of dollars to build these data centers and simply to train a single model. One version of a single model is about \$100 million just in computing power. No university grant can ever do that. Nobody, except these very limited numbers of incredibly wealthy corporations, is able to do it. Now, the Chinese have managed to do it somewhat cheaper, reportedly, but even so it’s private capital. So, we’ve got to kind of understand the global economy of these things and know that this is a very different beast from the Internet.\\
    There are thousands of people fine-tuning. All over the world, all over the country, they’re paid to fine-tune it and they change it, take out the bias, put in a bias. You know, focus on this, don’t focus on that. It’s all human control and human-generated. As soon as we are human, we are educators. What do we do? How do we intervene to offer guardrails and some control?\\
    Before you go into that, there is one other step that Mary says is kind of human-created. It’s the human inheritance on a number of levels. One is the Internet. When they scrape the Internet, when they scrape digitized text, they do decide to leak things out. They leak out private and personal information, for example, which is readily available on the Internet. They’re not worried about copyright, they steal copyright. That’s not copyright. They steal privacy, they steal everybody’s copyright. That’s where people have copyright. What they do is how you build the corpus. It’s like five billion words or something. It’s a defined amount of material, what they do, and then the training process about that works. It is not this process of working out the ways, the relationships of words with each other, but then, as Mary said, they do the fine-tuning. The fine-tuning has actually been relatively benign to me quite recently, so far, but the dangerous thing is it doesn’t have to be.\\
    Now, why we are going to say it’s been benign is because it’s basically a liberal California view of the world. Basically, the theology of the fine-tuning. What they do is they, you know, thousands of people, not just in the U.S. but around the world, peace workers, go in there and they ask questions. In addition, when you say tell me how to build, I’m a terrorist, tell me how to build a dirty nuclear bomb for the purpose of terrorism, they won’t answer, but they could answer. They’d just be told not to by the reinforcement. When you say, look, it’s the easiest way of suicide, they’re not going to give you the answer to that either. What happened was that sectors of the liberal elite in California, as well as a significant part of the United States political landscape, in contrast to the Republican Party and conservative groups, influenced the development of artificial intelligence systems. As a result, by incorporating values associated with a more liberal perspective, AI tends to present multiple viewpoints and avoid explicitly sexist or racist expressions. It’s a certain kind of California liberalism, and a very interesting footnote, by the way, is that now even the Chinese LLMs are having their fine-tuning done in California. Anyhow, what’s really frightening about this is that a malevolent political force of some kind could easily get in and tune it another way. Meanwhile, it’s very easy to turn around entirely the ideology of the Chinese LLMs. As easily as they tuned it to their ideology, another person could tune it to a different ideology.\\
    Before I get on to what we do, I (Mary) want to say that it also learns by the way we engage with it. We provide content; the content is already there. Every time we like something, every time we put in an email, we are contributors to that content market usage, because it’s learning from our bodies, you know, we watch it. No, that’s not generated AI, it’s actually learning from its prompts, but it collects data. The bigger picture is companies like Google, particularly, not OpenAI. OpenAI can get information from prompts, but Google has huge amounts of data from your emails, from your search history, so the AI is learning from our human experience all the time.\\
    Also, Google seems to be more sustainable going into the future than the other. So what do we do? That was the question. What have we done? Because we’ll talk about that also. We don’t have to send our learners into the wild. When using AI in a classroom, any kind of classroom, we have to curate the corpus. It is possible to curate a corpus by saying, I’m an English teacher, I’m going to put in these books, these tables, these images, these videos, creating a corpus, and our learners go into that corpus, not into the wild. Firstly, we have to be creators of content in a boundless corpus. Second, we need to create what we call rubric agents. How will our students access that corpus, and how will they be assessed against the goals that the course sets? So, those are two things we have to consider. Thirdly, we have to become what’s called prompt engineers, creating a set of prompts that goes with the corpus and a set of prompts that prepares our learners to engage with the knowledge that’s there. In this way, we put a barrier between the wild and what we want; anybody can do it. You can do it about music, about history, about any topic you want. You can’t do it without software. Why don’t you have the software? You have access to software in order to do it.\\
    For example, when it comes to prompts in AI, the fact that the system knows all your previous prompts means it’s getting to know you and giving answers which align with your interests and sensibilities because it’s got to know you. What we’ve done, for example, is we’ve been building this environment called Cyber Scholar, and what happens is everything the students do, the system records, and it gets to know the students. Thus, we’ve used exactly that same technique to do it. Now, we aren’t worried about privacy in the sense that we have a confined system. It’s completely closed, but we do go out to the foundation models for additional support. That goes in a way which is so highly de-identified. No names of the students are there, and not even the work they’re doing is there. It’s just tiny little packets going out all the time to foundation models. We build a separation from that. So, we’re not giving any student identities or data to the foundation models. Yet, the most important thing is the layer we build above the foundation model. We’ve got the foundation model there and what we’ve built is we’re not committed to any one foundation model. Also, we built a system where we can access any of all the foundation models, so, you know, check GPT, Cod, DeepSeek, the Chinese one, Mistral, the French one, all of them.\\
    Rubric agents are one of the conditions of knowing and what I want learning to be in this environment. This is a series of agents which interrogate students’ work from different perspectives. For example, if you take our Learning by Design Multiliteracy Schema, there’s a bit that deals with experiencing the known, a bit that deals with experiencing experiential learning, a bit that deals with concepts, a bit that deals with critical thinking. Each of these we configure as an agent, which is a kind of persona that will separately, and one persona at a time, one agent at a time, evaluate the student work and give them feedback on the basis of the rubric agents. Now, we call these rubric agents because they’re reminiscent of rubrics in education. In fact, it’s the criteria for learning the elements of what we consider to be successful learning, but also, it’s the elements of this environment. The second thing we have is an area called a knowledge ratio. I’ll mention what this is. This is for retrieval or measurement generation. Therefore, what we do is consider this to be a little bit like the textbook, except it can be huge. You can upload any resources the teachers curated: readings, references, textbooks, anything. I upload all of that, and that gets what’s called vectorized.\\
    There was a series of waves, so it supplements the donation model and you build kind of a corpus of your own. Now, both these things are pretty simple. It’s a series of forms where you write out your rubric agents. These are the things I want to be looking for in student learning, and these are the resources I want to use for this subject. The subject might be 18th-century poetry, English poetry. It might be cell biology. Whatever subject it is, you’re doing that. Then, what we can do is build this environment where we’ve heavily curated it to align with what learners are doing here. It’s an AI-based environment. That’s a quick overview of what we’ve been doing, anyhow, to experiment with these new technologies. In a way, the most important point is to get back to this idea that we’ve built this filter over the foundation models. The companies have been able to do it and then above that, we’ve built this specifically educational layer. One of the aspects of California liberalism is the customer’s always right. So, they’ll be nice to you, and they’ll patronize you and they’ll make you feel good. This is the customer service industry. For example, I’m a student and it says: how can I help you? Oh, I’ll write the essay for you. Right. Well, that’s not helpful. What we’ve done is we’ve built, and they’re the system meta-prompts created by the AI, the generative AI provider. Thus, we’ve built a place where we’ve said, don’t give answers, ask questions, make suggestions. In a sense, the AI becomes kind of a teacher, which keeps prompting the students to figure things out themselves.\\
    Now, I (Mary) started in a way to tell you about the theory of practice. We are not software designers, but in this moment, if we don’t become software designers, if we don’t collaborate with others, we can’t intervene. We’ve had to collaborate with others, you know, differently in order to bring our educational values to bear on AI and this digital moment. It’s really important because we’ve never had to do that before as educators. We wrote textbooks, the syllabi, papers, and then we lectured on our paper. Whereas now we have to have new tools as educators. This is really important because the tools that are being given to us by the commercial world, right, are mostly about deductive pedagogy that moves, you know, click, click, click, and teachers use it for behavior, controlling behavior. Thus, every child has an app for spelling and an app for mathematics—an app, right, repetitive, not collaborative. These are the commercial products that come back to us because we as educators are saying: what do we do? We have to collaborate to make this shift into creating what we’ve just demonstrated. That is possible because we’re in an American university, where we had access to resources, grant money, and where we can say it is possible to tame the AI, to harness it, and use it for ethical and productive reasons. It’s possible to do it, but it’s cyber-social. It’s the relationship between the human and the machine. We, as humans, have to calibrate that and not wait for the outside world to keep giving us these apps.\\
    Now, the learning management systems that are there, that we all use, and that our universities put millions of dollars on, are not architected to address the sorts of things we want as educators. Yet, that’s another barrier to moving into designing the tools that we need. It’s difficult, what we’re sharing with you, but it has a practical example, and it has lots of data that’s coming out now that says our students are not cheating. Our students are creative; they’re not cognitively offloading. All the dangers that we say are potentially there when you use AI, we’ve been able to address with the design that we’ve put into this Cyber Scholar platform, which is a pedagogical platform.
}
    \begin{description}
        \interviewcontent
    \end{description}
% Interview End
% Interview Begin
\setinterviewee{Mary Kalantzis and Bill Cope}{}{University of Illinois, Urbana- Champaign, Illinois, USA}{john@example.com}
\setinterviewer{Francis}{}{Universidade Federal de Minas Gerais, Colégio Técnico, Belo Horizonte, MG, Brasil}{victor.lim@nie.edu.sg}
\renewcommand{\interviewdate}{\today}
\renewcommand{\interviewcontext}{
    \textbf{Context:} This is a template interview.
}
\setspeaker{Qu}[\interviewerfname]
\setspeaker{An}[\intervieweefname]
\renewcommand{\interviewcontent}{
    \Qu You gave us many good examples and types for the design. That’s great! It is just a comment, not a question. But I’ve read the news about the Elon Musk’s experience in El Salvador. He has a relationship with El Salvador’s president. It’s a controversial government there. And they have the authority to implement Elon Musk’s AI. The name is GROK. I used it once. It’s not good AI. It’s a controversial AI, but the plan is to implement GROK in all schools in El Salvador. That’s 5,000 schools and 1 million students in El Salvador who will start using GROK, but they don’t have a project. They didn’t have a talk with the teachers in the schools, just using the AI in a school for the students, a little experience in a small country. It’s not a good experience. We have similar experiments like that in some states in Brazil. It’s not a government decision or a national decision, it’s a local decision and a policy based on votes because people have the experience: “I think it’s a good experience because AI technology is the future in schools”, but without a project, like you said, about the good experience with Cyber Scholar, it won’t work.
    \An GROK is fine-tuned by Elon Musk in existing ways and in particular ways. He’s not into Californian liberalism. He’s into a different kind of rival and it’s about ideologically influencing people. They’ve done that before, but we as educators cannot go into the wild. Also, no government accepts going into the wild. We have to create, as I said, curated spaces where our learners go in order to use the affordances of the new digital. We, as educators, need to create the conditions for how they experience it and how they’re able to move through the door. \\
    Actually, in our scope, you can go straight into any of the models without using a corpus, without using replication. You can. So that’s an option as well. Meanwhile, what we want to do is kind of, as researchers, we want to leave that question wide open. If there’s an image between GROK and GROK, we want to know what it is. We want to test what it is. We want to see the outcomes. Also, the other thing that we’re doing now, which is unique, is we’re tracking how much is created by the student and how much they take from AI. So you can track that. We have added that because that’s another thing we need to know. We need to know how much they’ve been independently creating and how much they access the AI. Now, we value the AI because it can be their partner, because it has all the knowledge. It has everything, but our students need to understand the relationship between what they want to say, what they need to do, and what the AI offers. That’s what we’ve done with the cyber scholar track problem that we’ve produced.
}
    \begin{description}
        \interviewcontent
    \end{description}
% Interview End


\section{Question 3}\label{sec-q3}
% Interview Begin
\setinterviewee{Mary Kalantzis and Bill Cope}{}{Organisation}{john@example.com}
\setinterviewer{Fei Victor Lim}{}{Organisation}{jane@example.com}
\renewcommand{\interviewdate}{\today}
\renewcommand{\interviewcontext}{
    \textbf{Context:} This is a template interview.
}
\setspeaker{Qu}[\interviewerfname]
\setspeaker{An}[\intervieweefname]
\renewcommand{\interviewcontent}{
    \Qu You have also argued for the end of assessment as we know it. With GenAI now capable of generating essays and other conventional learning products, relying on such artefacts as evidence of learning becomes increasingly untenable. How might GenAI support a move away from assessing only the end products of learning (artefacts created) towards capturing and valuing the processes of learning itself?
    \An We have always, in all our work, been concerned with assessment as being a negative thing that happens at the end of an experience that can be good or bad. It’s a judgmental thing, it doesn’t advance education, and it’s about short-term memory for the exam. We’ve always believed that assessment has to be formative. It has to happen as you go, assessment as you go, and learners need feedback at the point at which they need it rather than at the end, when it’s too late. In all our work, we’ve been arguing for a rethinking of assessment in education away from political reasons.\\
    Thus, what we’ve been able to do through the use of the digital is track performance and provide feedback. We’ve done it with peers. That’s one way, peer feedback, because peers learn from each other and it’s really important to evaluate each other’s work on the basis of a rubric. We have a rubric, so the peers respond. The teacher can respond too, but now we’ve added AI. The learner has three types of responses to their work: \textbf{peer review, AI review, and examiner or teacher review}. That’s built into this pedagogy platform that we have, and every learner has an extra plot. It’s a plot, right, that shows how they grow in learning, how they grow in the way they engage, and for us, how they grow in helping each other, or learning design. That is the quality of how they help each other with feedback, and now we’ve added how much AI they use. Those four things can all be tracked and you get collaborative learning, social learning, and just-in-time feedback. Also, the student decides what grade they want because they can see the grade. That’s motivational for them because our students all want to get better marks. They know they won’t get better marks because they’re seeing what knowledge they’re producing, what feedback, the quality of the feedback, how much AI, and levels of engagement. In this way, we’re going to track different things now, particularly in the online space. When we’re working online, you have to engage, you have to, if you value collaboration, give it a grade. If you value engagement, you have to give it a grade. So, we’ve figured out a system of providing grades that the learner can control themselves and advance in the way they want. The teacher can see all the grades, too. The teacher can see every extra plot that a student has, including the class one, and intervene as they like. But for us, we think and hope you’ll understand that the teacher’s role is as important as ever, but it’s different. We’re not going to have a future without teachers by just going to AI; we’ll just be brainwashing. We’ll still need educators who mediate the relationship between the human student and the machine knowledge.\\
    There will be one big change, though. Firstly, the relatively small change is that teachers have to put their learning resources into these knowledge bases; they have to do a rubric on it. The big change is that once that’s there, the AI becomes like an intermediary. It’s really the teacher that’s taught when the AI gets feedback. Thus, the teacher becomes this person with knowledgeable voices. Instead of it being a one-to-one classroom, it’s a one-to-one classroom besides the teacher, who is one-to-one with everybody simultaneously. That’s a really big and exciting possibility. The other exciting possibility is that at the school end, this is very cheap. This is just a web app, you know; we have to pay the token charges for the foundation model. There are some costs in this, and it’s impossible, even though they’re getting the private data to get away with the cost, but the costs are actually very modest compared to the historical costs and the labor costs of teachers, the historical costs of school buildings, the cost of textbooks. In other words, what we’ve got is the possibility of a really very different kind of learning universe, and also one at a very low cost. So, we are kind of both; we’re optimistic about this, but there are terrifying aspects of this as well. \\
    The other terrifying aspect for educators generally is that this is going to lead to the automation of much knowledge work. If you take the whole history of work in the modern era, mechanization and agriculture meant that very few people, far fewer people worked on farms. Mechanization of factories and robotics meant that far fewer people worked in factories. Also, the labor-intensive, high-education areas, like software development, accounting, and lawyers, are all going to be substantially automated. One of the real crisis points is the automation of knowledge work—labor-intensive, high-education knowledge work. Actually, low-education knowledge work too, like working in a call center, for example. One of the examples is the possibility of a massive unemployment crisis and the divide between the rich and the poor getting just worse and worse. One of the things is that it’s socially potentially cataclysmic. We have to have a politics ready beside us to deal with that, and the politics might be the politics of guaranteed income. It has to be a politics of redistribution of wealth. We’ve got to have a way to deal with that as well, outside of education.\\
    As educators, we need to show what’s possible. Even in our own university, what we’re doing is hard, right? This is a world-class, first-year, top university because millions of dollars have gone into the learning management system. It’s very hard to shift. Also, people are saying: oh! Let’s go back to face-to-face teaching or let’s just go back to tests with pen and paper. We’ve put every bit of our energy and every bit of our scholarship into saying ‘let’s create the tools that show that it’s possible to maintain our values and maintain the role of the educator, which is so important, in the education of future generations’. So that’s what we’ve done with cyber scholar. It’s not cyber org, it’s cyber scholar.\\
    In a sense, what we do is we go back to sovereignty to the rise of writing. What writing means is that we lost a whole lot of things about oral cultures. We’re from Australia, which had before colonisation these incredibly rich oral cultures, these rich communities. The loss of those things is a product of literacy. It destroys mental capacity to our memory; mental capacity is around. One of the things is we are in a world now where the world of literacy has produced a lot of good things—science, progress, and enlightenment—but a lot of bad things as well, such as bureaucracy, control, and differential power. Yet, both good and bad, it changed our minds. It changed human minds. Our sense is that this change now is going to be another huge change. Then, we will increasingly come to rely on AIs to help us think, to help us write. You know, literacy is a way of thinking, it’s a way of communicating, and it’s a profoundly important tool. If AI was with us all the time, it was going to change the way we do this, it would be our minds in the process. So, you know, our sense is that this is potentially a very big transition now that means: what do we do as educators? We have to be training students all the time on how to work with AI, use it productively and use its potential. So, we argue that that’s an enormous demand.\\
    Once the printing press and writing produced gatekeepers, questions arose: What is economical knowledge? Who gets published? Who gets to write? This moment allows more diversity potentially, more capacity for different voices to participate in meaning-making and sharing. We see that there’s a positive expanding of the production of knowledge if we train our people to be able to do it effectively and powerfully. Thus, we have a responsibility in that direction also.\\
    The other thing is multilingualism. I’m going (Bill) to make a provocative comment now about multilingualism. You won’t need to learn imperial languages anymore. You’ll be able to work in very small languages. You’ll have your ear pods in your ear and you’ll be able to speak. One of the great contradictions people think is that the language models are dominated by English; it’s actually the opposite. The opposite is that, in fact, they’re working with semantics of a number of languages and the ones in particular that are trained on multiple models. I’ll give you one example: a couple of languages in Rwanda, which have a few million speakers but almost no digital corpus. They’re just oral languages. Well, in fact, you can create a digital corpus with those instantly based on the fact that it’s leveraging the underlying semantics of multiple languages. What’s interesting is that the way in which tokenization happens, which is the basic mechanism in LLMs, is effectively working with semantics. It aligns very closely to a variant by notion of grammar, by the way. It’s effectively working with human grammars in general. So, what happens is that even languages that don’t have a big digital corpus can be very powerfully used in schools and for education, and even if you’re a learner of a very small, a minor language, you don’t need to learn any other language. Everything will be simultaneously translated. \\
    Now, this language teaching and learning is one of the areas that I think is going to be devastated by the technologies. Already, what we’ve seen is interpreting and translating have already been devastated. So, you know, there used to be 10 major programs in the U.S. Now there are four. I don’t know the exact numbers. That’s roughly what’s happened. They’re closing down very quickly. No human, no matter how good they are, can be as good as the AI in doing those things. Anyway, we’re trying to answer your questions with a narrative behind it, with a story behind it about our purpose. You know, what do we do as researchers and educators? What’s our role? Why do we do these things? So that’s why we responded to your questions in this bigger way.
}
    \begin{description}
        \interviewcontent
    \end{description}
% Interview End
% Interview Begin
\setinterviewee{Mary Kalantzis and Bill Cope}{}{Organisation}{john@example.com}
\setinterviewer{Vânia Castro}{}{Organisation}{jane@example.com}
\renewcommand{\interviewdate}{\today}
\renewcommand{\interviewcontext}{
    \textbf{Context:} This is a template interview.
}
\setspeaker{Qu}[\interviewerfname]
\setspeaker{An}[\intervieweefname]
\renewcommand{\interviewcontent}{
    \Qu Question 3 is a very interesting and complex question. One of the goals scholars have been striving for in education is that students should be less passive learners and more active knowledge makers. Students should be producers of multimodal artifacts, and teachers should assess these learning products and the way students have cited all the sources. However, now that GenAI can produce any multimodal text, there is a misconception  that learners should no longer produce them. \\
    When we think of evidence of learning, we can think of demonstrating knowledge to the teacher, but also for themselves. If educators continue to encourage students to produce artifacts, it means that students will have extraordinary learning opportunities, if it is implemented in ways that foster students’ creativity and curiosity. According to \textcite{freire1996}, in his book \textit{Pedagogy of Autonomy}, educators must reinforce the critical capacity of the learners. He states that curiosity, as a questioning restlessness, is an integral part of the vital phenomenon of learning. He also emphasizes that creativity cannot exist without the curiosity that drives us when we do and learn something. Teachers could create assessment activities that require creativity, in a way that students can add their own personal experiences and identity, layers that AI could not supply. \\
    Teachers should focus on both learning and assessment as a process, not as a final product. This idea is not new; \textcite{hyland2004,hyland2007} used to discuss the notion of writing as a process, but now we should expand that to other competences. \\
    When students produce artifacts, they develop numerous metacognition skills, such as knowledge organization and communication competence. They become aware of their strengths and weaknesses as learners, writers, and readers, express themselves, and register their ideas through publishing their artifacts not only for the teachers but for a real audience, classmates, and others.\\
    However, when assessment centers only on the final product and students know that the artifact they reproduce will be the “final test”, it may kill the process of creation/creativity; they will produce that only for the purpose of passing in school. This is what we need to rethink and change: “the assessment as we know it”. So, teachers should perceive the production of artifacts as a cognitive process.\\
    Students have spent too much time answering the right questions. It is more than time to teach them how to ask critical questions rather than give the right answers. It is time to reconsider the way we assess students’ artifacts, as there is a difference between learning activities, which could be, for example, any writing assignment, and assessment activities. For that, educators need some additional structures to measure learning, during and after the artifact creation. These include oral reasoning, AI justification and how students integrated AI tools into the creation process, sharing the prompt and the conversation students had with the LLMs, explanation of their decisions, and other strategies. This is all about academic integrity. Students should be transparent on how they use AI, but teachers also need to be role models and share how they have been using these tools, so that they can learn from each other and develop a relationship of trust. Students should learn some honest, fair, and ethical ways of using AI as tools, not replacing their capabilities.\\
    It is important to remember that teachers should not rely on a single type of assessment. If testing is the only measurement of learning, it disadvantages some students, especially the ones who experience anxiety on the day of the test, and even though they have enough knowledge to succeed. This is one of the ways to eliminate existing inequalities in classrooms. Equity values and practices should be evident in teachers’ practices.\\
    The whole idea is to focus on learning and assessment as a process, not as a final product, and making sure that learning comes first with critical thinking, and digital/AI technologies come second.
}
    \begin{description}
        \interviewcontent
    \end{description}
% Interview End


\section{Question 4}\label{sec-q4}
% Interview Begin
\setinterviewee{Vânia Castro}{}{Organisation}{john@example.com}
\setinterviewer{Camila Cárdenas}{}{Organisation}{jane@example.com}
\renewcommand{\interviewdate}{\today}
\renewcommand{\interviewcontext}{
    \textbf{Context:} This is a template interview.
}
\setspeaker{Qu}[\interviewerfname]
\setspeaker{An}[\intervieweefname]
\renewcommand{\interviewcontent}{
   \Qu We know that using AI applications carries risks and challenges that go beyond the formal educational setting, raising relevant ethical issues regarding the promotion of critical thinking and responsible citizenship. For instance, there are concerning effects related to the confirmation of cognitive biases among adolescents and young people who turn to chatbots for psychological guidance or to reinforce polarized political views, besides authorship and plagiarism concerns. What advice would you give teachers to deal with these issues and to adapt their pedagogical practices accordingly?
   \An There are many layers and dimensions for this question. One key suggestion is that teachers should invest in developing their AI literacy skills. It is crucial that they understand the concept of Artificial Intelligence (AI) and Generative AI (GenAI). They need to have a repertoire of AI tools and understand how to use them, but not superficially or in a way that the AI will replace their efforts. Having a variety of AI tools, developing a deep understanding of how these tools work, and knowing the principles behind them is essential for any educator. Teachers need to develop some critical thinking about AI tools and be aware of both benefits and limitations, which include bias, hallucination, data dependency, cognitive atrophy, algorithmic judgement, environmental impact, intellectual property, and many others. When teachers are aware of all those challenges, they become much more prepared to make decisions when using AI for education and to prepare their students to be critical users of these new tools. These limitations require thoughtful and in-depth considerations. \\
   When teachers address these issues with students, it is also important to bring examples to the classroom that connect to students’ reality. For example, they could show students that by filling 500ml of water and then pouring it out, the wasted water represents the hidden cost each time they generate a question to Large Language Models (LLMs), such as ChatGPT, Copilot, or others. According to \posscite{heikkila2023} archive from MIT Technology Review, writing a 100-word email with AI uses more energy than four LED bulbs lit for an hour.\\
   Bias is another challenge teachers should examine carefully. The Data Provenance Initiative report, published in 2024, analyzed the public dataset in LLMs, from 1990 to 2024, and only 0.2\% of this data come from Africa and South America, including Brazil. So, this illustrates how Brazil is underrepresented in the data from these models. Based on that, we could reflect critically: what is the level of biased or incomplete content generated by these GenAI tools? To understand more about the underrepresentation and the issues related to algorithms in LLMs, I recommend reading the works by Júlio Araújo, a Brazilian expert in algorithmic racism in AI systems. Some questions that teachers should keep in mind: how can I prepare my students to critically evaluate any content generated from AI? What steps should we take to ensure that our students will identify bias, incomplete content, or hallucinations? If teachers invest time in developing their AI literacy skills, they will be better equipped to integrate them, strategically, into their pedagogical practices.
}
    \begin{description}
        \interviewcontent
    \end{description}
% Interview End



\section{Question 5}\label{sec-q5}
% Interview Begin
\setinterviewee{Vânia Castro}{}{Organisation}{john@example.com}
\setinterviewer{Fei Victor Lim}{}{Organisation}{jane@example.com}
\renewcommand{\interviewdate}{\today}
\renewcommand{\interviewcontext}{
    \textbf{Context:} This is a template interview.
}
\setspeaker{Qu}[\interviewerfname]
\setspeaker{An}[\intervieweefname]
\renewcommand{\interviewcontent}{
   \Qu In your work, you note that we are all becoming “cyborgs” because of our increasing reliance on digital technologies in everyday literacy practices. How does this shift reshape our understanding of multiliteracies and multimodal literacy?
   \An Firstly, it is worth recalling the main aspects of multiliteracies theory so that we can align them with the metaphor “cyborgs”. Broadly speaking, the multiliteracies notion takes into account the diversity present in educational settings, multimodality, and the use of digital technologies aligned with pedagogy. Another important aspect is the idea that teachers should prepare students to become democratic citizens and embrace diverse, cross-cutting themes in teaching practices. Multiliteracies also came to supplement traditional reading and writing, and expand literacy practices by recognizing the distinct ways in which people engage in contemporary communication, as recalled by \textcite{kalantzis2015}. \\
   The “cyborgs” metaphor is alighted to the idea of “cognitive prostheses” by \textcite{kalantzis2020}, which they define as material artifacts to support our representations. In this perspective, digital tools become part of human cognitive systems, extending human thinking and meaning-making possibilities. Multimodal literacy involves interaction with all kinds of technologies, including artificial intelligence, not as something separate from human activity, but as elements that reshape how we make meaning, read, write, and communicate.
   \\We need to recognize that multiliteracies are plural and that multimodality is central, so literacy pedagogy must change. Meaning should be constructed through combinations of modes, and learners should be active knowledge designers and citizens, making meaning within an expanded landscape intertwined with digital and emerging tools.
}
    \begin{description}
        \interviewcontent
    \end{description}
% Interview End

\printbibliography\label{sec-bib}
% if the text is not in Portuguese, it might be necessary to use the code below instead to print the correct ABNT abbreviations [s.n.], [s.l.]
%\begin{portuguese}
%\printbibliography[title={Bibliography}]
%\end{portuguese}


%full list: conceptualization,datacuration,formalanalysis,funding,investigation,methodology,projadm,resources,software,supervision,validation,visualization,writing,review
\begin{contributors}[sec-contributors]
\authorcontribution{Mary Kalantzis}[conceptualization,investigation,validation]
\authorcontribution{Bill Cope}[conceptualization,investigation,validation]
\authorcontribution{Vania Carvalho de Castro}[conceptualization,investigation,validation]
\authorcontribution{Záira Bomfante dos Santos}[supervision,visualization,writing,review]
\authorcontribution{Francis Arthuso Paiva}[supervision,visualization,writing,review]
\authorcontribution{Fei Victor Lim}[supervision,visualization,writing,review]
\authorcontribution{Camila Cárdenas Neira}[investigation,methodology,resources]
\end{contributors}


\begin{dataavailability}
\txtdataavailability{nodata} % options: dataavailable, dataonly, databody, datanotav, nodata
\end{dataavailability}


\end{document}

